\chapter{操作系统简史}
\label{app:os-history}

今天的电脑和手机上，跑的多半是这四种操作系统之一：Windows、Linux、macOS 和 Android。它们都是“现代操作系统”：都能同时运行很多个程序，都用虚拟内存把程序隔开（第~\ref{chap:memory}~章），都通过系统调用为程序提供服务（第~\ref{chap:os}~章）。从内核的角度看，它们做的事情大同小异。

可是用起来，它们给人的感觉却很不一样。这个附录先讲讲它们是从哪里来的，再说说为什么会有这种“不一样”的感觉。这里讲的东西，不懂也完全不影响读这本书；要真正弄明白一个操作系统是怎样工作的，往往要花上好几年。现在只要知道“原来是这么回事”就够了。

\section{一张家谱}

操作系统的历史，像一棵枝繁叶茂的家谱树（图~\ref{fig:os-tree}）。今天这四个“表亲”，其实只来自两个老祖宗：一个是 1969 年诞生在贝尔实验室的 \emph{Unix}，另一个是 1981 年随 IBM PC 一起出世的 \emph{DOS}。

\begin{figure}[htbp]
  \centering
  \begin{tikzpicture}[
      os/.style={draw, rounded corners=3pt, font=\small, minimum height=6mm, inner sep=3pt, align=center},
      unix/.style={os, fill=green!15},
      ms/.style={os, fill=blue!12},
      apple/.style={os, fill=gray!15},
      link/.style={-{Stealth}, thick, gray!70!black},
      idea/.style={-{Stealth}, dashed, gray},
      yr/.style={font=\scriptsize, gray!70!black, anchor=east},
      line width=0.6pt,
    ]
    % 年份刻度
    \foreach \y/\t in {0/1969, -1.3/1980 年代, -2.6/1990 年代, -3.9/2000 年代, -5.2/今天}
      \node[yr] at (-0.4,\y) {\t};
    \draw[gray!40] (-0.3,0.3) -- (-0.3,-5.5);

    % Unix 一支
    \node[unix] (unix)  at (2.4,0)    {Unix\\\scriptsize 贝尔实验室};
    \node[unix] (bsd)   at (1.2,-1.3) {BSD\\\scriptsize 伯克利};
    \node[unix] (minix) at (3.6,-1.3) {Minix\\\scriptsize 教学用};
    \node[unix] (linux) at (3.6,-2.6) {Linux\\\scriptsize 1991};
    \node[apple](next)  at (1.2,-2.6) {NeXTSTEP};
    \node[apple](mac)   at (1.2,-3.9) {Mac OS X\\\scriptsize 2001};
    \node[apple](macos) at (0.4,-5.2) {macOS};
    \node[apple](ios)   at (1.5,-5.2) {iOS};
    \node[unix] (distro)at (2.85,-5.2) {Ubuntu 等\\\scriptsize 发行版};
    \node[unix] (andr)  at (5.0,-5.2) {Android\\\scriptsize 2008};

    \draw[link] (unix) -- (bsd);
    \draw[idea] (unix) -- (minix);
    \draw[idea] (minix) -- node[right, font=\scriptsize, gray] {启发} (linux);
    \draw[link] (bsd) -- (next);
    \draw[link] (next) -- (mac);
    \draw[link] (mac) -- (macos);
    \draw[link] (mac) -- (ios);
    \draw[link] (linux) -- (distro);
    \draw[link] (linux) -- (andr);

    % 微软一支
    \node[ms] (dos)   at (6.6,-1.3) {MS-DOS\\\scriptsize 1981};
    \node[ms] (win9x) at (6.2,-2.6) {Windows 3.x\\95 / 98};
    \node[ms] (nt)    at (9.0,-2.6) {Windows NT\\\scriptsize 1993};
    \node[ms] (xp)    at (7.65,-3.9) {Windows XP\\\scriptsize 2001};
    \node[ms] (win)   at (7.65,-5.2) {Windows 10 / 11};
    \draw[link] (dos) -- (win9x);
    \draw[link] (win9x) -- (xp);
    \draw[link] (nt) -- (xp);
    \draw[link] (xp) -- (win);

    % 图例
    \node[unix, font=\scriptsize] at (5.4,0.15) {Unix 一家};
    \node[apple, font=\scriptsize] at (6.6,0.15) {苹果};
    \node[ms, font=\scriptsize] at (7.6,0.15) {微软};
  \end{tikzpicture}
  \caption{操作系统的家谱（大大简化了）}
  \label{fig:os-tree}
\end{figure}

\paragraph{Unix 一家。}
1969 年，贝尔实验室的 Ken Thompson 和 Dennis Ritchie 写出了 Unix。为了把 Unix 移植到别的机器上，Ritchie 发明了 C 语言，后来 Unix 的大部分代码都改用 C 重写——这就是 C 语言和操作系统那么“亲”的原因。Unix 后来分出了许多支，其中最有名的是加州大学伯克利分校的 \emph{BSD}\lazynote{BSD}{Berkeley Software Distribution，加州大学伯克利分校发布的 Unix 版本。}。前言里提到的滑动窗口，就是从 4.4BSD 的代码里学来的。

\paragraph{Linux。}
1991 年，芬兰大学生 Linus Torvalds 受教学用的小系统 Minix 启发，在自己的 80386 电脑上写了一个内核，取名 Linux（第~\ref{sec:bits}~节）。它不是 Unix 的后代，却照着 Unix 的样子来写，所以用起来几乎一样。Linux 本身只是一个内核，加上 GNU 的工具（gcc、glibc、bash……）和各种软件，打包在一起，才成了 Ubuntu、Debian 这样的“发行版”\lazynote{发行版}{Linux 内核加上工具、软件和安装程序，打包成一个完整的系统，比如 Ubuntu。}。今天，世界上绝大多数的服务器、全部前 500 名的超级计算机，跑的都是 Linux。

\paragraph{苹果。}
苹果的 Mac 电脑在 1984 年就有了图形界面，但它今天的 macOS 却另有来历：乔布斯离开苹果后创办的 NeXT 公司，做了一个基于 BSD 的系统 NeXTSTEP；1997 年苹果买下 NeXT，乔布斯回归，NeXTSTEP 就变成了 2001 年的 Mac OS X，也就是今天的 macOS。iPhone 上的 iOS，是它的兄弟。所以，macOS 其实是一个正宗的 Unix。

\paragraph{微软。}
1981 年，IBM PC 用的操作系统是微软的 MS-DOS：一个只有黑底白字命令行、一次只能运行一个程序的小系统。后来的 Windows 3.x、95、98，都是在 DOS 的基础上加了图形界面。与此同时，微软请来了 Dave Cutler，从头写了一个全新的、真正现代的内核：Windows NT（1993）。2001 年的 Windows XP 起，家用的 Windows 也换成了 NT 内核，一直用到今天的 Windows 11。第~\ref{sec:win32}~节里见过的 \texttt{ntdll.dll}，名字里的 NT 就是这么来的。

\paragraph{Android。}
2008 年，谷歌推出了手机操作系统 Android。它的内核就是 Linux，上面却换上了自己的一整套东西（连 C 标准库都换成了自己的 Bionic，第~\ref{sec:library}~节）。所以世界上数量最多的 Linux，其实装在大家的手机里。

\begin{figure}[htbp]
  \centering
  \begin{tikzpicture}[line width=0.8pt]
    % 企鹅（Linux）
    \begin{scope}[shift={(0,0)}]
      \fill[black] (0,1.1) ellipse (0.62 and 0.95);
      \fill[white] (0,0.95) ellipse (0.42 and 0.72);
      \fill[black] (0,2.15) circle (0.45);
      \fill[white] (-0.15,2.22) circle (0.12) (0.15,2.22) circle (0.12);
      \fill[black] (-0.13,2.22) circle (0.05) (0.13,2.22) circle (0.05);
      \fill[orange] (-0.15,2.0) -- (0.15,2.0) -- (0,1.86) -- cycle;
      \fill[orange] (-0.4,0.12) ellipse (0.25 and 0.1) (0.4,0.12) ellipse (0.25 and 0.1);
      \node[font=\small] at (0,-0.35) {Linux};
    \end{scope}
    % 窗户（Windows）
    \begin{scope}[shift={(2.6,0.45)}]
      \fill[red!70]    (-0.65,0.85) rectangle (-0.05,1.45);
      \fill[green!60!black] (0.05,0.85) rectangle (0.65,1.45);
      \fill[blue!60]   (-0.65,0.15) rectangle (-0.05,0.75);
      \fill[yellow!80!orange] (0.05,0.15) rectangle (0.65,0.75);
      \node[font=\small] at (0,-0.8) {Windows};
    \end{scope}
    % 苹果（macOS）
    \begin{scope}[shift={(5.0,0.25)}]
      \fill[gray!60] (0,0.9) circle (0.62);
      \fill[white] (0.62,1.05) circle (0.22);
      \fill[gray!60] (0.08,1.62) ellipse (0.12 and 0.25);
      \node[font=\small] at (0,-0.6) {macOS};
    \end{scope}
    % 机器人（Android）
    \begin{scope}[shift={(7.4,0)}]
      \fill[green!55!black] (-0.5,0.4) rectangle (0.5,1.4);
      \fill[green!55!black] (-0.5,1.5) arc[start angle=180, end angle=0, radius=0.5];
      \fill[white] (-0.2,1.75) circle (0.06) (0.2,1.75) circle (0.06);
      \draw[green!55!black, line width=2pt, line cap=round] (-0.25,1.95) -- (-0.4,2.2) (0.25,1.95) -- (0.4,2.2);
      \draw[green!55!black, line width=4pt, line cap=round] (-0.68,0.6) -- (-0.68,1.25) (0.68,0.6) -- (0.68,1.25)
                                                          (-0.22,0.4) -- (-0.22,0.1) (0.22,0.4) -- (0.22,0.1);
      \node[font=\small] at (0,-0.35) {Android};
    \end{scope}
    \node[draw, rounded corners, fill=yellow!15, font=\small, align=center] at (3.7,3.1)
      {我们其实是亲戚！\\（除了那个窗户……它也学了不少）};
  \end{tikzpicture}
  \caption{今天的“四大天王”}
  \label{fig:four-os}
\end{figure}

\section{为什么感觉不一样}

既然内核做的事大同小异，为什么换一个系统，就会觉得浑身不自在？原因很多：界面长得不一样，按键的习惯不一样，自带的软件不一样……其中最常让人摸不着头脑的一个，是\emph{文件是怎么组织的}。

\subsection{Windows：一个个盘}

在 Windows 上，每一块硬盘（或者硬盘上的每一个分区、每一个 U 盘）都是一个“盘”，用一个字母来叫：\texttt{C:}、\texttt{D:}、\texttt{E:}……每个盘都有自己的一棵文件夹树，互不相干：
\begin{minted}[linenos=false]{text}
C:\
+-- Windows\          Windows 自己
+-- Program Files\    安装的软件
`-- Users\
    `-- xiaoming\     你的文件：桌面、文档、下载……
D:\
`-- work\
    `-- mypaper\      这本书
\end{minted}
路径用反斜杠 \texttt{\textbackslash} 分隔，比如 \texttt{D:\textbackslash work\textbackslash mypaper}。为什么从 C 开始？因为 \texttt{A:} 和 \texttt{B:} 当年留给了软盘驱动器——那是 DOS 时代的事，可这个规矩一直留到了今天。大小写也不分：\texttt{HELLO.C} 和 \texttt{hello.c} 是同一个文件。

\subsection{Linux：只有一棵树}

Linux（以及 macOS、Android）里没有“盘”的概念。整个系统只有一棵树，树根就叫 \texttt{/}（根目录）。所有的东西都长在这棵树上：硬盘、U 盘、甚至键盘和屏幕，都只是树上的某个文件夹或某个文件。在 WSL 的终端里敲 \texttt{ls /}，看到的就是根目录下的这些文件夹：
\begin{minted}[linenos=false]{text}
$ ls /
bin   boot  dev  etc  home  lib  media  mnt  opt  proc
root  run   sbin  srv  sys  tmp  usr  var
\end{minted}
它们的名字都很短，是 Unix 时代留下来的缩写（表~\ref{tab:fhs}）。

\begin{table}[htbp]
  \begin{lazywide}
  \caption{Linux 根目录下常见的文件夹}
  \label{tab:fhs}
  \begin{tabular}{lL{3.8cm}L{6.0cm}}
    \toprule
    文件夹 & 原意 & 放的是什么 \\
    \midrule
    \texttt{/home} & home & 每个用户的“家”：\texttt{/home/xiaoming} 就相当于 Windows 的 \texttt{C:\textbackslash Users\textbackslash xiaoming} \\
    \texttt{/bin}、\texttt{/usr/bin} & binary（二进制） & 程序：\texttt{ls}、\texttt{gcc} 都在这里 \\
    \texttt{/lib}、\texttt{/usr/lib} & library & 函数库：\texttt{libc.so.6} 就在 \texttt{/usr/lib/x86\_64-linux-gnu/} 里 \\
    \texttt{/usr/include} & include & 头文件：\texttt{stdio.h} 就在这里（第~\ref{sec:stdio-h}~节） \\
    \texttt{/etc} & 原来是“et cetera”（其他） & 系统的设置文件 \\
    \texttt{/tmp} & temporary & 临时文件，重启就没了 \\
    \texttt{/dev} & device & 设备：硬盘、终端……都以文件的样子出现在这里 \\
    \texttt{/proc} & process & 正在运行的程序的信息，第~\ref{sec:layout}~节读过的 \texttt{/proc/self/maps} 就在这里 \\
    \texttt{/mnt}、\texttt{/media} & mount & 其他硬盘、U 盘“挂”在这里 \\
    \texttt{/root} & & 超级用户（管理员）的家，别和根目录 \texttt{/} 弄混 \\
    \bottomrule
  \end{tabular}
  \end{lazywide}
\end{table}

那另外一块硬盘去哪了？在 Linux 上，要先把它“挂载”（mount）\lazynote{挂载}{把一块硬盘（或 U 盘）接到目录树的某个文件夹上，从那里就能访问它。}到树上的某个文件夹，它才能被看到。WSL 就是这样做的：Windows 的 \texttt{C:} 盘挂在 \texttt{/mnt/c}，\texttt{D:} 盘挂在 \texttt{/mnt/d}。所以这本书在 Windows 里叫 \texttt{D:\textbackslash work\textbackslash mypaper}，在 WSL 里叫 \texttt{/mnt/d/work/mypaper}——是同一个文件夹，只是两个系统叫法不同。Linux 的路径用正斜杠 \texttt{/} 分隔，而且分大小写：\texttt{Hello.c} 和 \texttt{hello.c} 是两个不同的文件。

\begin{figure}[htbp]
  \centering
  \begin{tikzpicture}[line width=0.8pt]
    % 左：Windows 的一排柜子
    \foreach \x/\l in {0/C, 1.5/D, 3.0/E} {
      \filldraw[fill=brown!35] (\x,0) rectangle +(1.2,1.8);
      \foreach \y in {0.15,0.75,1.35} \draw (\x+0.1,\y) rectangle +(1.0,0.5);
      \foreach \y in {0.4,1.0,1.6} \fill (\x+0.6,\y) circle (0.04);
      \node[draw, fill=white, font=\bfseries] at (\x+0.6,2.15) {\l:};
    }
    \node[font=\small] at (2.1,-0.45) {Windows：一个盘一个柜子};
    % 右：Linux 的一棵树
    \begin{scope}[shift={(7.6,0)}]
      \fill[brown!60] (-0.15,0) rectangle (0.15,1.1);
      \draw[brown!60!black, line width=2pt] (0,0.9) -- (-0.9,1.6) (0,0.9) -- (0.9,1.6) (0,1.1) -- (0,1.9);
      \fill[green!45!black, opacity=0.85] (0,2.0) ellipse (1.7 and 0.95);
      \node[draw, fill=white, font=\ttfamily\small] at (0,0.45) {/};
      \foreach \x/\y/\t in {-1.0/1.9/home, -0.2/2.55/etc, 0.75/2.2/usr, -0.55/1.45/bin, 0.75/1.55/mnt} {
        \node[fill=yellow!60, font=\ttfamily\scriptsize, inner sep=1.5pt, rounded corners=1pt] at (\x,\y) {\t};
      }
      \node[font=\small] at (0,-0.45) {Linux：所有东西都长在一棵树上};
    \end{scope}
    % 一个从柜子走到树下的困惑小人
    \draw (5.2,1.35) circle (0.18);
    \draw (5.2,1.17) -- (5.2,0.55);
    \draw (5.2,0.55) -- (5.05,0) (5.2,0.55) -- (5.35,0);
    \draw (5.2,1.0) -- (4.95,1.25) -- (5.0,1.5);
    \draw (5.2,1.0) -- (5.45,0.8);
    \node[draw, rounded corners, fill=white, text=red!70!black, font=\small\bfseries] (q) at (5.2,3.3) {我的 D 盘呢？};
    \draw (q.south) -- (5.2,1.6);
    \draw[gray, line width=1pt] (-0.4,0) -- (9.8,0);
  \end{tikzpicture}
  \caption{柜子和树}
  \label{fig:drives-tree}
\end{figure}

两种办法都能用，谈不上谁对谁错。一开始觉得别扭，只是因为不习惯。不必一下子全弄明白，先接受“它们就是这样组织的”，用着用着就熟了。

\section{为什么在 \texorpdfstring{\NoCaseChange{Windows}}{Windows} 上也用 \texorpdfstring{\NoCaseChange{GCC}}{GCC}}

这本书从头到尾都用 \texttt{gcc} 编译程序，连 Windows 的程序也是用 MinGW 版的 gcc 编译的（第~\ref{sec:win32}~节）。为什么不用微软自己的编译器？

最简单的原因：\textbf{它是免费的}，而且是自由软件\lazynote{自由软件}{可以自由使用、研究、修改和分发的软件，源代码是公开的。}，谁都可以下载、使用、甚至修改它的源代码。GCC 是 GNU 项目在 1987 年推出的，最早叫“GNU C Compiler”，后来支持的语言越来越多，就改叫“GNU Compiler Collection”。今天，它能为几十种 CPU 生成代码：x86、ARM、RISC-V……等我们以后自己造出 CPU，用来给它编译程序的，也会是 gcc。

当然，好用的编译器不止这一个。微软的 Visual Studio（用的是 MSVC 编译器）有免费的社区版，在 Windows 上写图形界面的程序很方便；苹果的 Xcode 和 Linux 上也很流行的 Clang，同样是免费的。学 C 语言，用哪个都行。本书选 gcc，是因为它在 Linux、Windows（通过 WSL 或 MinGW）上都能用，命令也都一样。

\section{终端和图形界面}

操作系统有两张脸：一张是黑底白字的\emph{终端}\lazynote{终端}{输入命令、显示文字结果的窗口。里面运行的 \texttt{bash} 叫做 shell。}（命令行），一张是有窗口、图标、鼠标的\emph{图形界面}（GUI）。很多人以为终端是“老古董”、图形界面才“先进”，或者反过来，觉得会敲命令才是“高手”。其实两者没有好坏之分，只是适合做不同的事（图~\ref{fig:cli-gui}）。

\begin{figure}[htbp]
  \centering
  \begin{tikzpicture}[line width=0.8pt]
    % 终端
    \filldraw[fill=black!85, draw=gray] (0,0) rectangle (4.2,2.6);
    \fill[gray!50] (0,2.35) rectangle (4.2,2.6);
    \foreach \x/\c in {0.15/red!70, 0.35/yellow!80!orange, 0.55/green!60!black} \fill[\c] (\x,2.475) circle (0.06);
    \node[anchor=north west, align=left, font=\ttfamily\tiny, text=green!70] at (0.1,2.25)
      {\$ gcc hello.c -o hello\\\$ ./hello\\Hello, world!\\\$ make test\\全部通过\\\$ \_};
    \node[font=\small] at (2.1,-0.35) {终端：敲命令};
    % 图形界面
    \begin{scope}[shift={(5.4,0)}]
      \filldraw[fill=cyan!15, draw=gray] (0,0) rectangle (4.2,2.6);
      \fill[gray!25] (0,0) rectangle (4.2,0.3);
      \foreach \x/\y/\t in {0.4/1.9/文档, 1.2/1.9/图片, 2.0/1.9/音乐} {
        \fill[yellow!70!orange] (\x-0.25,\y-0.2) rectangle (\x+0.25,\y+0.2);
        \node[font=\tiny] at (\x,\y-0.38) {\t};
      }
      \filldraw[fill=white, draw=gray] (2.3,0.55) rectangle (3.95,1.75);
      \fill[blue!50] (2.3,1.55) rectangle (3.95,1.75);
      \node[font=\tiny] at (3.12,1.1) {画图};
      \draw[fill=white] (3.3,0.75) -- (3.3,1.1) -- (3.4,1.0) -- (3.5,1.2) -- (3.56,1.17) -- (3.46,0.98) -- (3.58,0.97) -- cycle;
      \node[font=\small] at (2.1,-0.35) {图形界面：点鼠标};
    \end{scope}
    % 猫和狗
    \node[draw, rounded corners, fill=yellow!15, font=\small] at (2.1,3.15) {狗：敲命令才是真本事！};
    \node[draw, rounded corners, fill=yellow!15, font=\small] at (7.5,3.15) {猫：我用爪子点一下就行。};
  \end{tikzpicture}
  \caption{终端和图形界面：做的是不同的事，没有谁比谁高级}
  \label{fig:cli-gui}
\end{figure}

\paragraph{终端更适合的事：}
\begin{itemize}
  \item \textbf{重复和批量。}把一百张照片改名，在终端里是一行命令，用鼠标得点一百次。本书里的 \texttt{make test}，一下子就把几十个程序都跑了一遍。
  \item \textbf{精确和可重复。}命令可以写进文件里，下次原样再跑一遍，结果一模一样。书中每一段程序输出，都是这样自动生成的。
  \item \textbf{远程管理。}服务器往往没有显示器，管理员通过网络连上去，靠的全是终端。
  \item \textbf{编程。}编译、调试、查看系统调用（\texttt{strace}）……程序员的大部分工具，都是在终端里用的。
\end{itemize}

\paragraph{图形界面更适合的事：}
\begin{itemize}
  \item \textbf{看和画。}看照片、看视频、修图、画画，没有屏幕和鼠标简直没法做。
  \item \textbf{探索。}第一次用一个软件，菜单和按钮会告诉你它能做什么，不用先背命令。
  \item \textbf{日常使用。}上网、聊天、写文档、玩游戏——大多数人大多数时候，都在用图形界面。
\end{itemize}

其实，这本书就是两张脸一起用写出来的：在终端里用 Vim 写稿、用 \texttt{latexmk} 编译，再在图形界面的 PDF 阅读器里看排版的效果。而且，不管是终端还是图形界面，它们都只是普通的程序：终端里的 \texttt{bash}、Windows 的桌面 \texttt{explorer.exe}，最后都要通过系统调用去请内核办事。对内核来说，它们没有任何区别。

\section{小结}

\begin{itemize}
  \item 今天的 Windows、Linux、macOS、Android 都是现代操作系统，做的事大同小异。它们分别来自 Unix 和 DOS 两个老祖宗：macOS 来自 BSD，Android 用的是 Linux 内核，今天的 Windows 用的是 NT 内核。
  \item 让人觉得“不一样”的原因之一，是文件的组织方式：Windows 用 \texttt{C:}、\texttt{D:} 这样的盘；Linux 只有一棵以 \texttt{/} 为根的树，其他硬盘挂载到树上。
  \item GCC 免费、自由，能为很多种 CPU 编译程序，在 Linux 和 Windows 上都能用。
  \item 终端和图形界面是操作系统的两张脸，各有擅长的事，没有好坏之分。
  \item 这些不必一下子全弄懂。操作系统要学好几年，先知道“原来是这样”就好。
\end{itemize}
